\documentclass[12pt,letterpaper]{article}

\usepackage[utf8]{inputenc}
\usepackage[T1]{fontenc}
\usepackage[spanish]{babel}
\usepackage{geometry}
\usepackage{hyperref}
\usepackage{xcolor}
\usepackage{enumitem}
\usepackage{longtable}
\usepackage{array}
\usepackage{listings}

\geometry{margin=2.5cm}
\hypersetup{
  colorlinks=true,
  linkcolor=blue,
  urlcolor=blue
}

\lstdefinestyle{comandos}{
  basicstyle=\ttfamily\small,
  breaklines=true,
  frame=single,
  columns=fullflexible
}

\title{Guía de despliegue de CapyCode}
\author{Documento técnico básico}
\date{\today}

\begin{document}

\maketitle

\section{Introducción}

Este documento explica cómo instalar y ejecutar el proyecto CapyCode en una computadora o servidor. La idea es dejar una guía práctica para una persona que no necesariamente sabe mucho de servidores, pero que necesita levantar el sistema de forma ordenada.

CapyCode no es solamente una página estática. No funciona como una página simple de Netlify donde solo se suben archivos HTML. Este proyecto usa:

\begin{itemize}
  \item un servidor en Node.js;
  \item una base de datos MySQL;
  \item archivos web en HTML, CSS y JavaScript;
  \item modelos 3D cargados en el navegador;
  \item variables privadas en un archivo \texttt{.env}.
\end{itemize}

Por eso, para que la página funcione, el servidor de Node.js y la base de datos deben estar funcionando al mismo tiempo.

\section{Resumen rápido de ejecución}

La forma general de levantar el proyecto es esta:

\begin{enumerate}
  \item Instalar Node.js y npm.
  \item Instalar MySQL.
  \item Crear/configurar el archivo \texttt{.env}.
  \item Instalar las dependencias una sola vez con:
\end{enumerate}

\begin{lstlisting}[style=comandos]
npm install
\end{lstlisting}

\begin{enumerate}[resume]
  \item Cargar la base de datos una sola vez con:
\end{enumerate}

\begin{lstlisting}[style=comandos]
sudo mysql -u root -p < schema.sql
\end{lstlisting}

En Windows normalmente no se usa \texttt{sudo}; el comando equivalente suele ser:

\begin{lstlisting}[style=comandos]
mysql -u root -p < schema.sql
\end{lstlisting}

\begin{enumerate}[resume]
  \item Ejecutar la página con:
\end{enumerate}

\begin{lstlisting}[style=comandos]
npm start
\end{lstlisting}

\section{Dependencias necesarias en Windows 11}

Asumiendo una instalación normal o ``de fábrica'' de Windows 11, normalmente se necesita instalar lo siguiente.

\subsection{Node.js y npm}

Node.js es el entorno que ejecuta el servidor del proyecto. npm es el gestor de paquetes que descarga las librerías necesarias.

\begin{itemize}
  \item Instalar Node.js desde \url{https://nodejs.org/}.
  \item Se recomienda usar una versión LTS moderna, por ejemplo Node.js 20 o 22.
  \item npm viene incluido al instalar Node.js.
\end{itemize}

Para comprobar que quedó instalado:

\begin{lstlisting}[style=comandos]
node -v
npm -v
\end{lstlisting}

\subsection{MySQL}

MySQL es la base de datos donde se guardan usuarios, progreso, skins, grupos, preguntas respondidas y otros datos del sistema.

\begin{itemize}
  \item En Windows se puede instalar MySQL Community Server.
  \item También puede usarse una instalación de MySQL en otro servidor.
  \item Lo importante es que el servidor MySQL esté encendido mientras CapyCode está corriendo.
\end{itemize}

Para probar conexión:

\begin{lstlisting}[style=comandos]
mysql -u root -p
\end{lstlisting}

\subsection{Git}

Git no es obligatorio para ejecutar el proyecto si ya se tiene la carpeta, pero sí es muy útil para descargarlo desde GitHub o mantener versiones.

\begin{lstlisting}[style=comandos]
git --version
\end{lstlisting}

\section{Dependencias del proyecto}

Las dependencias JavaScript se instalan con:

\begin{lstlisting}[style=comandos]
npm install
\end{lstlisting}

Este comando lee \texttt{package.json} y \texttt{package-lock.json}. No hace falta correrlo cada vez que se inicia la página. Se corre una vez al instalar, o cuando cambian las dependencias.

\begin{longtable}{|>{\raggedright\arraybackslash}p{3.2cm}|>{\raggedright\arraybackslash}p{3cm}|>{\raggedright\arraybackslash}p{7cm}|}
\hline
\textbf{Paquete} & \textbf{Versión instalada} & \textbf{Uso dentro del proyecto} \\
\hline
\endhead
\texttt{express} & \texttt{4.22.2} & Servidor web de Node.js. Sirve la página, archivos estáticos y rutas API. \\
\hline
\texttt{three} & \texttt{0.172.0} & Motor/librería 3D usada en el navegador para mostrar mapas, personajes y portales. \\
\hline
\texttt{mysql2} & \texttt{3.22.4} & Conector para que Node.js pueda hablar con MySQL usando promesas. \\
\hline
\texttt{dotenv} & \texttt{16.6.1} & Carga variables privadas desde el archivo \texttt{.env}. \\
\hline
\texttt{bcryptjs} & \texttt{2.4.3} & Sirve para manejar contraseñas de forma segura mediante hash. \\
\hline
\texttt{jsonwebtoken} & \texttt{9.0.3} & Maneja tokens de sesión para login y autenticación. \\
\hline
\end{longtable}

\subsection{Compatibilidad aproximada}

El proyecto usa dependencias modernas, pero no demasiado raras. En general:

\begin{itemize}
  \item Node.js 20 o 22 debería funcionar bien.
  \item Node.js 18 probablemente también funcione, pero es mejor usar LTS reciente.
  \item Express 4.x es compatible con el estilo actual del servidor.
  \item Three.js puede cambiar bastante entre versiones, así que conviene quedarse con \texttt{0.172.0} salvo que se pruebe bien.
  \item MySQL 8 es una opción recomendada. MySQL 5.7 o MariaDB pueden funcionar en varios casos, pero no conviene asumirlo sin probar el esquema.
  \item \texttt{package-lock.json} ayuda a instalar las mismas versiones exactas.
\end{itemize}

\section{Configuración del archivo .env}

El archivo \texttt{.env} es necesario porque guarda datos privados o configuraciones que no deben escribirse directamente en el código.

Ejemplo de estructura:

\begin{lstlisting}[style=comandos]
PORT=3000
DB_HOST=localhost
DB_USER=capycode
DB_PASSWORD=tu_password
DB_NAME=capycode
JWT_SECRET=cambia_esto_por_un_texto_largo
OPENAI_API_KEY=opcional
OPENAI_MODEL=gpt-4o-mini
\end{lstlisting}

\subsection{Por qué no se debe subir .env a GitHub}

El archivo \texttt{.env} puede contener contraseñas de la base de datos y llaves de API. Por eso debe estar en \texttt{.gitignore}. Si se sube por accidente, otra persona podría usar esas credenciales.

En producción, si se usa un hosting como Render, Railway u otro, normalmente estas variables se escriben en el panel del hosting, no en un archivo físico.

\section{Preparar la base de datos}

Antes de ejecutar el proyecto, hay que crear las tablas de MySQL. Este proyecto incluye el archivo:

\begin{lstlisting}[style=comandos]
schema.sql
\end{lstlisting}

En Linux o en un servidor donde se use \texttt{sudo}, se puede correr:

\begin{lstlisting}[style=comandos]
sudo mysql -u root -p < schema.sql
\end{lstlisting}

En Windows, normalmente sería:

\begin{lstlisting}[style=comandos]
mysql -u root -p < schema.sql
\end{lstlisting}

Este paso se hace una sola vez al preparar la base de datos. Después de eso, MySQL debe quedar corriendo mientras la página esté activa.

\section{Instalación y ejecución}

\subsection{Instalación inicial}

Dentro de la carpeta del proyecto:

\begin{lstlisting}[style=comandos]
npm install
\end{lstlisting}

\subsection{Ejecución normal}

Para iniciar el servidor:

\begin{lstlisting}[style=comandos]
npm start
\end{lstlisting}

El proyecto usa el script definido en \texttt{package.json}:

\begin{lstlisting}[style=comandos]
"start": "node server.js"
\end{lstlisting}

Por defecto, si \texttt{PORT} no está definido, el servidor usa el puerto \texttt{3000}. Entonces se abre en:

\begin{lstlisting}[style=comandos]
http://localhost:3000
\end{lstlisting}

\section{Notas para hosting}

En un hosting real, no basta con subir los archivos como una página HTML normal. El hosting debe permitir ejecutar Node.js y conectarse a una base de datos MySQL.

El servidor que aloje el proyecto debe tener:

\begin{itemize}
  \item Node.js instalado o soportado por el proveedor.
  \item Las variables de entorno configuradas.
  \item Acceso a una base de datos MySQL.
  \item El comando de instalación: \texttt{npm install}.
  \item El comando de inicio: \texttt{npm start}.
\end{itemize}

Si la base de datos está en otro servidor, entonces \texttt{DB\_HOST}, \texttt{DB\_USER}, \texttt{DB\_PASSWORD} y \texttt{DB\_NAME} deben apuntar a ese servidor.

\section{Problemas comunes}

\subsection{La página no abre}

Revisar:

\begin{itemize}
  \item que \texttt{npm start} no haya mostrado errores;
  \item que el puerto sea correcto;
  \item que se esté abriendo \texttt{http://localhost:3000} o el dominio correcto.
\end{itemize}

\subsection{Error de conexión a MySQL}

Posibles causas:

\begin{itemize}
  \item MySQL no está encendido.
  \item El usuario o contraseña en \texttt{.env} están mal.
  \item La base de datos no existe.
  \item No se ejecutó \texttt{schema.sql}.
\end{itemize}

\subsection{Login no funciona}

Puede ser problema de base de datos o de sesión. Revisar:

\begin{itemize}
  \item que existan usuarios en la base de datos;
  \item que \texttt{JWT\_SECRET} esté definido;
  \item que el servidor no esté mostrando errores en consola.
\end{itemize}

\subsection{Los modelos 3D no cargan}

Revisar:

\begin{itemize}
  \item que existan los archivos en \texttt{public/assets/objects};
  \item que el navegador no muestre errores 404 en la consola;
  \item que la versión de \texttt{three} instalada sea la esperada.
\end{itemize}

\subsection{Python AI Help no aparece}

El botón de ayuda con IA solo aparece si existe una llave API configurada. Revisar:

\begin{lstlisting}[style=comandos]
OPENAI_API_KEY=tu_llave
\end{lstlisting}

Si no se configura, el botón no aparece y el chat queda deshabilitado.

\section{Conclusión}

El despliegue de CapyCode requiere más pasos que una página estática, porque combina servidor, base de datos y frontend 3D. La parte más importante es instalar dependencias con \texttt{npm install}, preparar MySQL con \texttt{schema.sql}, configurar bien el archivo \texttt{.env} y ejecutar el servidor con \texttt{npm start}.

Si esos puntos están bien, el proyecto debería poder ejecutarse de forma estable en una computadora local o en un servidor que soporte Node.js y MySQL.

\end{document}
